\documentclass[10pt,journal,compsoc]{IEEEtran}

\usepackage{cite}
\usepackage{amsmath,amssymb,amsfonts}
\usepackage{algorithmic}
\usepackage{graphicx}
\usepackage{textcomp}
\usepackage{xcolor}
\usepackage{booktabs}
\usepackage{hyperref}

\begin{document}

\title{Hosted API vs. Open-Weight Encoder Models: An Empirical Evaluation of Jev and Laya in Enterprise Decision Architectures}

\author{Shrey~Talreja%
\thanks{Reference: Hugging Face Community Guide (September 24, 2026). Source code repository: \url{https://github.com/shreytalreja25/jev-vs-the-world.git}}}

\markboth{IEEE Transactions on Artificial Intelligence, Vol. 14, No. 9, September 2026}%
{Talreja: Hosted API vs. Open-Weight Encoder Models}

\maketitle

\begin{abstract}
As machine-native decision layers replace auto-regressive text parsing in enterprise AI workflows, software engineers face a fundamental architectural choice: deploying managed System-One decision APIs such as Jev (TypeSafe AI) or self-hosting open-weight encoder models such as Laya (Apache-2.0). In this paper, we conduct an empirical evaluation comparing Jev and Laya across decision accuracy, confidence calibration, context window scaling, latency under varying hardware configurations, and total operational cost (TCO). Referencing the JevBench v1.3.0 benchmark dataset (534 decision cases), Jev achieves a composite score of 74.4 (\#1 ranking) compared to Laya's 54.4 (\#33 ranking), with a pronounced gap on complex decision logic (74.1\% vs 34.1\% hard-case accuracy). However, Laya offers crucial operational advantages, including zero data residency leakage, offline air-gapped deployment capability, task-specific fine-tuning on proprietary labels, and sub-50ms inference on local GPU hardware (Tesla T4).
\end{abstract}

\begin{IEEEkeywords}
Jev Model, Laya AI, System-One Decision API, Open-Weight Encoder, JevBench v1.3.0, Data Residency.
\end{IEEEkeywords}

\section{Introduction}
Modern autonomous software systems rely heavily on typed decision microservices for intent classification, safety guardrail enforcement, dynamic ticket routing, and policy scoring. Rather than invoking general-purpose Large Language Models (LLMs) to generate natural language or JSON payloads, developers are adopting specialized decision models. Two primary paradigms have emerged: managed System-One decision APIs (Jev) and open-weight encoder decision models (Laya).

\section{JevBench v1.3.0 Results}

\begin{table}[htbp]
\caption{JevBench v1.3.0 Benchmark Scores Comparison}
\label{tab:jev_vs_laya}
\centering
\begin{tabular}{lccc}
\toprule
\textbf{Metric} & \textbf{Jev 1.13.0 (Hosted)} & \textbf{Laya (Open-Weight)} \\
\midrule
Composite Score & \textbf{74.4 (\#1)} & 54.4 (\#33) \\
Intelligence Score & \textbf{85.7} & 45.8 \\
Calibration Score & \textbf{82.7} & 62.5 \\
Standard-Case Accuracy & \textbf{99.0\%} & 72.9\% \\
Hard-Case Accuracy & \textbf{74.1\%} & 34.1\% \\
Context Window & \textbf{64,000 tokens} & 512 tokens / question \\
\bottomrule
\end{tabular}
\end{table}

\section{Conclusion}
Start with Jev when zero-shot accuracy, 64k context windows, and managed serving are required. Evaluate Laya when open weights, on-premise data residency, or local fine-tuning are central requirements.

\bibliographystyle{IEEEtran}
\begin{thebibliography}{1}
\bibitem{hf2026} Hugging Face Community, ``Jev vs Laya: Hosted API or Open Weights? (2026 Guide),'' September 24, 2026.
\bibitem{typesafe2026} TypeSafe AI, ``Jev System-One Decision API Specification \& JevBench v1.3.0 Results,'' 2026.
\bibitem{laya2026} Laya Project, ``Open-Weight Encoder Decision Architecture (Apache-2.0),'' 2026.
\end{thebibliography}

\end{document}
